\section{A general entropy bound from a reversible observation}
\label{sec:general-bound}

The exact four-state calculation admits an observation-only extension. A
negative off-diagonal entry of the principal logarithm forces every
compatible generator to use a nonprincipal logarithm branch. The resulting
oscillatory component has a discrete minimum size. Combining this fact
with stationary probability currents gives the following explicit bound.
This section keeps the state space fixed; the hidden-state problem is
treated separately.

\subsection{Setting and theorem}

Let $P$ be an $n\times n$ stochastic matrix, $n\geq3$, reversible with
respect to a strictly positive probability vector $p$, and suppose every
eigenvalue of $P$ is strictly positive. Fix an observation interval
$\tau>0$ and put
\[
 D=\operatorname{diag}(p),\qquad
 R=D^{1/2}PD^{-1/2},\qquad
 H=\tau^{-1}\log R.
\]
Thus $R$ is symmetric positive definite and $H$ is symmetric. Throughout
this section $\|\cdot\|_F$ denotes the Frobenius norm. Define
\begin{equation}
 T=-\frac{\log\det P}{\tau},\qquad
 d^2=\sum_{i\ne j}[-H_{ij}]_+^2,\qquad
 p_{\min}=\min_i p_i,
 \label{eq:gb-data}
\end{equation}
where $[x]_+=\max\{x,0\}$. Let
$p_{(1)}\leq p_{(2)}\leq\cdots\leq p_{(n)}$ denote the sorted
stationary probabilities, and write $s=p_{(1)}+p_{(2)}$.

A compatible generator is a real matrix $Q$ satisfying
\[
 Q\mathbf1=0,\qquad q_{ij}\geq0\quad(i\ne j),\qquad
 e^{\tau Q}=P.
\]
The entropy-production convention is
\begin{equation}
 \sigma(Q)=\sum_{i<j}
 (p_iq_{ij}-p_jq_{ji})
 \log\frac{p_iq_{ij}}{p_jq_{ji}}.
 \label{eq:gb-entropy}
\end{equation}
An edge with both fluxes zero contributes zero; one with exactly one
positive flux contributes $+\infty$. Natural logarithms are used, and
every state is even under the ordinary time reversal. Equation
\eqref{eq:gb-entropy} is a statistical irreversibility rate. Its
thermodynamic interpretation requires the usual physical assumptions
relating transition-rate ratios to entropy exchange.

\begin{theorem}[Entropy forced by a nonprincipal embedding]
\label{thm:general-bound}
Suppose $d>0$ and at least one compatible generator exists. Then $T>0$,
$p$ is stationary for every compatible generator, and
\begin{equation}
 \boxed{\displaystyle
 \sigma(Q)\geq
 \frac{\max\{12p_{\min},\,2s\}}{T}
 \left(d^2+\frac{8\pi^2}{\tau^2}\right)
 }
 \label{eq:gb-main}
\end{equation}
for every compatible $Q$. The compatible-generator set is compact, and
$\sigma$ attains its extended-valued minimum on this set. In particular,
the minimum is strictly positive, although it can be $+\infty$.

For the uniform stationary distribution, the bound simplifies to
\begin{equation}
 \sigma(Q)\geq\frac{12}{nT}
 \left(d^2+\frac{8\pi^2}{\tau^2}\right).
 \label{eq:gb-uniform}
\end{equation}
\end{theorem}

The bound is not asserted to equal the exact entropy minimum. Its role
is to give a closed-form positive certificate even when repeated
observed eigenvalues produce continuous families of compatible
logarithms. The proof is given in full in the next four subsections.

\subsection{Stationarity, the total exit rate, and attainment}

The communication classes of $Q$ and $e^{\tau Q}$ agree. One way to see
this is uniformization: for sufficiently large $c>0$,
\[
 e^{\tau Q}=e^{-c\tau}
 \sum_{k=0}^{\infty}\frac{(c\tau)^k}{k!}
 \left(I+\frac{Q}{c}\right)^k.
\]
The summands are nonnegative, and an entry is positive precisely when
the corresponding states are connected by a directed path of positive
rates, allowing the path of length zero.

Reversibility and $p>0$ imply that $P$, after a simultaneous permutation
of rows and columns, is block diagonal over irreducible communicating
classes. There are no transient classes. Therefore every compatible
$Q$ has the same block structure and is irreducible on each block.
The unique stationary distribution of a generator on one such block is
invariant under its exponential, so it equals the normalized restriction
of $p$ to that block. Consequently $pQ=0$. This argument includes the
irreducible case and the possible one-state blocks.

The determinant identity gives
\[
 \det P=\det(e^{\tau Q})=e^{\tau\operatorname{tr}Q}.
\]
Since $\operatorname{tr}Q$ is real, it follows without any logarithm-branch
ambiguity that
\begin{equation}
 \sum_i(-q_{ii})=-\operatorname{tr}Q=T.
 \label{eq:gb-trace}
\end{equation}
Every exit rate and every off-diagonal rate is therefore bounded by
$T$. If $T=0$, all rates vanish and $P=I$, contradicting $d>0$.

The generator constraints together with $e^{\tau Q}=P$ form a closed
subset of a bounded finite-dimensional set. Thus the nonempty set of
compatible generators is compact. The function
\[
 (a,b)\longmapsto(a-b)\log(a/b)
\]
with the conventions in \eqref{eq:gb-entropy} is nonnegative and lower
semicontinuous on $[0,\infty)^2$. At $(0,0)$, nonnegativity gives lower
semicontinuity; at a point with one positive and one zero coordinate,
the function diverges to $+\infty$. Hence $\sigma$ is lower
semicontinuous and attains an extended-valued minimum on the compact
compatible-generator set.

\subsection{An exact identity for the logarithm branches}

For a compatible generator define
\[
 A=D^{1/2}QD^{-1/2},\qquad
 S=\frac{A+A^T}{2},\qquad K=\frac{A-A^T}{2},\qquad
 U=S-H,
\]
and set $B=A-H=U+K$. The matrices $S$ and $U$ are symmetric, and $K$
is skew-symmetric. Off-diagonal entries of $A$ are nonnegative, so
$S_{ij}\geq0$ for $i\ne j$. Therefore
\begin{equation}
 \|U\|_F^2\geq\sum_{i\ne j}[-H_{ij}]_+^2=d^2.
 \label{eq:gb-negative-part}
\end{equation}

The matrix $A$ commutes with $R=e^{\tau A}$. It consequently preserves
every eigenspace of the symmetric matrix $R$, and commutes with
$H=\tau^{-1}\log R$. Equivalently, $H$ can be expressed as a polynomial
in $R$ by interpolation on its distinct positive eigenvalues. It
follows that
\[
 e^{\tau B}=e^{\tau A}e^{-\tau H}=I.
\]
Every eigenvalue of $B$ is thus of the form $2\pi i k_\ell/\tau$,
where $k_\ell\in\mathbb Z$. Moreover, $B$ is diagonalizable over
$\mathbb C$. Indeed, a nonzero nilpotent part $N$ of a Jordan block
cannot exponentiate to the identity, because
\[
 e^{\tau N}-I=N\,h(N),\qquad h(0)=\tau\ne0,
\]
and $h(N)$ is invertible. Since $B$ is real, its nonreal eigenvalues
occur in conjugate pairs.

Taking traces, without assuming that $B$ is orthogonally diagonalizable,
gives
\[
 \operatorname{tr}(B^2)
 =-\frac{4\pi^2}{\tau^2}\sum_{\ell=1}^n k_\ell^2.
\]
On the other hand, symmetry and skew symmetry give
$\operatorname{tr}(UK)=0$,
$\operatorname{tr}(U^2)=\|U\|_F^2$, and
$\operatorname{tr}(K^2)=-\|K\|_F^2$. Thus
\begin{equation}
 \boxed{\displaystyle
 \|K\|_F^2=\|U\|_F^2+
 \frac{4\pi^2}{\tau^2}\sum_{\ell=1}^n k_\ell^2.
 }
 \label{eq:gb-branch-identity}
\end{equation}

If all $k_\ell$ were zero, diagonalizability would imply $B=0$.
Then $A=H$, contradicting the negative off-diagonal entry of $H$.
Therefore a nonzero conjugate pair occurs and
$\sum_\ell k_\ell^2\geq2$. Combining
\eqref{eq:gb-negative-part} and \eqref{eq:gb-branch-identity} yields
\begin{equation}
 \|K\|_F^2\geq d^2+\frac{8\pi^2}{\tau^2}.
 \label{eq:gb-skew-floor}
\end{equation}

The term $8\pi^2/\tau^2$ is an exact branch obstruction, not a
small-perturbation approximation. More generally, it applies to every
nonprincipal embedding of a reversible positive-spectrum $P$, even if
its principal logarithm also happens to be a generator.

This reasoning also explains why the present setting necessarily has
repeated observed eigenvalues. A nonzero branch and its conjugate
must lie in the same real eigenspace of $R$. With simple positive
spectrum, every real logarithm is the principal logarithm.

\subsection{An elementary entropy estimate}

For an unordered pair $i<j$, write
\[
 a_{ij}=p_iq_{ij},\qquad b_{ij}=p_jq_{ji},\qquad
 c_{ij}=a_{ij}+b_{ij},\qquad j_{ij}=a_{ij}-b_{ij}.
\]
Here $c_{ij}$ is the total edge traffic and $j_{ij}$ is its signed
stationary current. The scalar inequality
\begin{equation}
 (a-b)\log\frac ab\geq\frac{2(a-b)^2}{a+b}
 \label{eq:gb-scalar}
\end{equation}
holds with the stated limiting conventions. For $a,b>0$, it follows
from $x\operatorname{atanh}x\geq x^2$, with
$x=(a-b)/(a+b)$. Edges with $c_{ij}=0$ contribute zero on both sides.

Define
\[
 \kappa=\max_{i<j}\frac{c_{ij}}{2p_ip_j}.
\]
Since $K_{ij}=j_{ij}/(2\sqrt{p_ip_j})$,
\begin{equation}
 \sigma(Q)\geq2\sum_{i<j}\frac{j_{ij}^2}{c_{ij}}
 \geq\frac{2}{\kappa}\|K\|_F^2.
 \label{eq:gb-kappa-bound}
\end{equation}
This entropy-to-skew estimate, including stronger nonlinear versions,
is established prior work; its relation to the literature is stated
below.

Let $x_i=-q_{ii}$. Stationarity makes the total incoming and outgoing
probability flux at state $i$ equal to $p_ix_i$. Each of $a_{ij}$ and
$b_{ij}$ is at most $p_ix_i$, and likewise at most $p_jx_j$.
Consequently
\[
 c_{ij}\leq2p_ix_i,\qquad c_{ij}\leq2p_jx_j,
\]
and
\[
 T\geq x_i+x_j\geq
 \frac{c_{ij}}{2p_i}+\frac{c_{ij}}{2p_j}.
\]
It follows that
\[
 \kappa\leq\max_{i<j}\frac{T}{p_i+p_j}=\frac Ts.
\]
Substitution into \eqref{eq:gb-kappa-bound} proves
\begin{equation}
 \sigma(Q)\geq\frac{2s}{T}\|K\|_F^2.
 \label{eq:gb-simple-constant}
\end{equation}

\subsection{A circulation improvement}

There is a stronger constant when the stationary probabilities are
uniform or not too uneven. We prove
\begin{equation}
 \sigma(Q)\geq\frac{12p_{\min}}{T}\|K\|_F^2.
 \label{eq:gb-circulation-constant}
\end{equation}
The case $\sigma(Q)=+\infty$ is immediate, and the case $K=0$ is
trivial, so assume finite entropy and a nonzero current.

Fix one orientation for every unordered pair $e=\{i,j\}$, and retain
the corresponding signed current $j_e$ and traffic $c_e$. Define
\[
 h_e=\frac{p_i+p_j}{2p_ip_j},\qquad
 L=\sum_e |j_e|\sqrt{h_e},\qquad
 N=\sum_e\frac{j_e^2}{p_ip_j}=2\|K\|_F^2.
\]
The current has zero divergence at every vertex. Expanding the sum of
all off-diagonal rates and using this fact gives the exact identity
\begin{align}
 T&=\sum_{e=\{i,j\}}
 \left[\frac{c_e+j_e}{2p_i}+\frac{c_e-j_e}{2p_j}\right]\notag\\
 &=\sum_e h_ec_e+
 \frac12\sum_{e=\{i,j\}}j_e\left(\frac1{p_i}-\frac1{p_j}\right)
 =\sum_e h_ec_e.
 \label{eq:gb-weighted-traffic}
\end{align}
The second sum vanishes because it is the pairing of the
divergence-free current with the gradient of $i\mapsto1/p_i$.
By Cauchy--Schwarz and \eqref{eq:gb-scalar},
\begin{equation}
 \sigma(Q)\geq2\sum_e\frac{j_e^2}{c_e}
 \geq\frac{2L^2}{T}.
 \label{eq:gb-lone-bound}
\end{equation}

Orient every nonzero-current edge in the direction of its current.
This nonnegative, divergence-free flow decomposes into directed simple
cycles with nonnegative amplitudes. For completeness, follow positive
edges until a vertex repeats, subtract the smallest flow on the
resulting cycle, and repeat. Each subtraction removes at least one
positive edge. No two-cycle occurs, because a nonzero net current
cannot have both orientations on the same unordered pair. Every
cycle therefore has length at least three.

For a simple cycle $C$, let $z_C$ denote its unit signed edge current
in the previously fixed orientation, and put
\[
 \ell_C=\sum_{e\in C}\sqrt{h_e},\qquad
 n_C=\sum_{\{i,j\}\in C}\frac1{p_ip_j},\qquad
 \rho=\min_C\frac{\ell_C^2}{n_C},
\]
where the minimum may be taken over all simple cycles of the complete
graph on the $n$ states. Write the decomposition as
$j=\sum_C\gamma_Cz_C$, with $\gamma_C\geq0$. It is sign-compatible,
so there is no cancellation on any edge and
\[
 L=\sum_C\gamma_C\ell_C,\qquad
 \frac jL=\sum_C\frac{\gamma_C\ell_C}{L}
                    \frac{z_C}{\ell_C}.
\]
The last expression is a convex combination. Convexity of the squared
weighted norm $x\mapsto\sum_e x_e^2/(p_ip_j)$ gives
\[
 \frac{N}{L^2}
 \leq\sum_C\frac{\gamma_C\ell_C}{L}\frac{n_C}{\ell_C^2}
 \leq\frac1\rho.
\]
Thus \eqref{eq:gb-lone-bound} yields the sharper, cycle-dependent
inequality
\begin{equation}
 \sigma(Q)\geq\frac{4\rho}{T}\|K\|_F^2.
 \label{eq:gb-cycle-bound}
\end{equation}

It remains to bound $\rho$ by observable stationary probabilities.
Set $a_i=p_{\min}/p_i\in(0,1]$. On any simple cycle,
\begin{align}
 \frac{\ell_C^2}{p_{\min}n_C}
 &=\frac{\left(\sum_{\{i,j\}\in C}
              \sqrt{(a_i+a_j)/2}\right)^2}
          {\sum_{\{i,j\}\in C}a_ia_j}\notag\\
 &\geq\frac{\left(\sum_{i\in C}a_i\right)^2}
              {\sum_{\{i,j\}\in C}a_ia_j}
 \geq3.
 \label{eq:gb-cycle-clique}
\end{align}
The first inequality uses $\sqrt{x}\geq x$ on $[0,1]$ and the fact
that every cycle vertex occurs twice in the edge sum.

Here is an elementary proof of the last inequality. Normalize the
nonnegative vertex weights so that their sum is one, and consider
$F=\sum_{\{i,j\}\in C}a_ia_j$. If two positive-weight vertices are
nonadjacent, hold their total weight fixed and move all of it to the
vertex with the larger weighted neighbor sum. Because the two
vertices are nonadjacent, $F$ is linear under this redistribution and
does not decrease. Repeating reduces the positive support to a
clique. A simple cycle has clique size $r\leq3$. On such a clique,
\[
 F=\frac{1-\sum_i a_i^2}{2}
 \leq\frac{1-1/r}{2}\leq\frac13.
\]
The original $F$ is no larger, proving
$(\sum_i a_i)^2\geq3\sum_{\{i,j\}\in C}a_ia_j$.
Consequently $\rho\geq3p_{\min}$, and
\eqref{eq:gb-cycle-bound} proves
\eqref{eq:gb-circulation-constant}.

For uniform $p_i=1/n$, one has $\ell_C^2/n_C=|C|/n$, so
$\rho=3/n$ when all edges are allowed. The constant $12/(nT)$ is
optimal for the entropy-to-skew inequality: take an almost-equilibrium
three-state cycle with the remaining states isolated, or connected
by arbitrarily weak reversible rates if irreducibility is required.
More explicitly, clockwise and counterclockwise rates $u+\varepsilon$
and $u-\varepsilon$ on a triangle give
\[
 \lim_{\varepsilon\to0}
 \frac{\sigma(Q)T}{\|K\|_F^2}=\frac{12}{n}.
\]
This establishes sharpness only of the entropy-to-skew constant, not
of the final bound \eqref{eq:gb-main}.

\begin{proof}[Completion of the proof of Theorem~\ref{thm:general-bound}]
The stationarity, compactness, and attainment statements were proved
above. Taking the stronger of \eqref{eq:gb-simple-constant} and
\eqref{eq:gb-circulation-constant}, then using
\eqref{eq:gb-skew-floor}, proves \eqref{eq:gb-main}. For uniform $p$,
$\max\{12p_{\min},2s\}=12/n$, giving
\eqref{eq:gb-uniform}.
\end{proof}

\subsection{The rational four-state example}

For the example with eigenvalues $1,2^{-8},2^{-18},2^{-18}$, uniform
stationary distribution, and observation interval one,
\[
 T=44\log2,\qquad H_{12}=H_{21}=-\frac{\log2}{2},\qquad
 d^2=\frac{(\log2)^2}{2}.
\]
Theorem~\ref{thm:general-bound} therefore supplies the explicit
certificate
\begin{equation}
 \sigma(Q)\geq
 \frac{3}{44\log2}
 \left(\frac{(\log2)^2}{2}+8\pi^2\right)
 =7.790264194785\ldots .
 \label{eq:gb-example}
\end{equation}
The separately solved four-state optimization gives a larger minimum,
$26.1789116783\ldots$. The difference records the information discarded
by a general inequality; there is no claim that
\eqref{eq:gb-example} recovers the family-specific optimum.

\subsection{Prior art and scope of the contribution}

The ingredients have established precedents. Jia's
\href{https://arxiv.org/abs/1605.03502}{reversible embedding theorem}
(2016) gives uniqueness and the principal-logarithm criterion for a
reversible embedding. Kolchinsky, Ohga, and Ito's
\href{https://arxiv.org/html/2304.01714v4}{thermodynamic spectral bounds}
(2024), particularly Appendix C.1, give
\eqref{eq:gb-kappa-bound} and stronger nonlinear estimates. Neither
the general matrix-logarithm machinery nor that entropy-to-skew
estimate is claimed as new here.

Bauer, Seifert, and van der Meer's
\href{https://arxiv.org/html/2412.15642v2}{stroboscopic inference paper}
enumerates compatible logarithms in the nondegenerate case. Section
3.3 notes the continuous families arising from repeated eigenvalues;
the estimator described in Section 4.1 is assigned zero in the
degenerate cases outside that enumeration. The present theorem
supplies a positive analytic bound for a specific degenerate regime:
reversible positive-spectrum observations whose principal logarithm
has a negative off-diagonal entry.

The observation-only combination of
\eqref{eq:gb-branch-identity}--\eqref{eq:gb-skew-floor} with the
determinant and stationary-probability constants, and the circulation
refinement above, were derived and checked in this investigation.
Their precise statements were not located in the sources reviewed.
That is a scoped originality assessment, not a certification of
historical priority.

The full observed state space and the fixed, homogeneous sampling
model are essential to \eqref{eq:gb-main}. The constants do not
automatically transfer to models with additional hidden states,
arbitrarily small hidden stationary probabilities, or unbounded
hidden rates. Dependence on stationary probabilities also separates
this theorem from proposed bounds using only an oscillatory generator
eigenvalue; Shiraishi's
\href{https://arxiv.org/html/2609.34352v1}{September 2026 counterexample}
rules out a broad class of the latter bounds. Finally, exact
reversibility and repeated eigenvalues should not be inferred from
noisy data without an uncertainty analysis. The theorem is a precise
inverse-problem statement; a statistical estimator requires
additional work.
